% ============================================================================
%  Section 6: Thực nghiệm và huấn luyện mô hình
% ============================================================================
\section{Thiết lập thực nghiệm}
\label{sec:experiments}

Chương này mô tả cấu hình được lưu trong repository, phạm vi dữ liệu của từng mô hình và quy trình tạo các artifact đánh giá. Những thông tin không còn log gốc để xác minh, như thời gian huấn luyện hoặc epoch dừng thực tế, không được trình bày như kết quả đo.

\subsection{Trạng thái prediction artifact}
\label{subsec:artifact-validity}

% Generated by scripts/generate_report_artifacts.py.
\begin{table}[htbp]
  \centering
  \caption{Trạng thái bằng chứng của các prediction artifact đã khóa bằng SHA-256.}
  \label{tab:artifact-validity}
  \small
  \begin{tabularx}{\textwidth}{lL{3cm}cL{2.7cm}X}
    \toprule
    \textbf{Artifact} & \textbf{Trạng thái} & \textbf{$n$} & \textbf{SHA-256} & \textbf{Phạm vi sử dụng} \\
    \midrule
Baseline & Hợp lệ & 234 & \code{a16c4d903e\ldots} & Bảng chính, RQ1, latency thành phần \\
VietOCR + \model{LayoutXLM} & Hợp lệ, kèm sensitivity check & 234 & \code{ea94f274d2\ldots} & Bảng chính, RQ1, latency thành phần \\
\model{Donut} & Không hợp lệ; chỉ chẩn đoán & 234 & \code{2cc970fd2a\ldots} & Phụ lục/chẩn đoán; không xếp hạng \\
    \bottomrule
  \end{tabularx}
\end{table}


Cả ba phương pháp Baseline, LayoutXLM và Donut đều có đủ 234 bản ghi dự đoán trên tập test, không có bản ghi nào bị lỗi (\code{status=error}) và đều được đưa vào so sánh chính. Với LayoutXLM, báo cáo bổ sung sensitivity check loại bốn ID có sai lệch tọa độ ảnh. Với Donut, lỗi cắt chuỗi 20 token trước đây đã được khắc phục hoàn toàn bằng cách đặt tham số sinh \code{generation\_max\_length=768}, giúp phục hồi đầy đủ năng lực suy luận của mô hình. SHA-256 trong bảng khóa đúng nội dung các artifact được phân tích.

% ============================================================================
\subsection{Môi trường và cấu hình}
\label{subsec:exp_env}

Các checkpoint được xây dựng cho môi trường GPU NVIDIA Tesla T4. Cấu hình phần mềm của dự án sử dụng Python, PyTorch, Hugging Face Transformers, PaddleOCR, VietOCR và RapidFuzz. Hạt giống chia dữ liệu được đặt là $42$. Do repository không lưu đầy đủ ảnh chụp môi trường và log package của lần huấn luyện tạo checkpoint, các phiên bản thư viện trong tài liệu hướng dẫn chỉ được xem là cấu hình dự kiến để tái lập, không phải bằng chứng về môi trường chạy lịch sử.

\begin{table}[htbp]
  \centering
  \caption{Cấu hình huấn luyện được khai báo cho Donut và LayoutXLM.}
  \label{tab:hyperparams}
  \begin{tabular}{lll}
    \toprule
    \textbf{Siêu tham số} & \textbf{\model{Donut}} & \textbf{\model{LayoutXLM}} \\
    \midrule
    Mô hình nền tảng & \code{donut-base} & \code{layoutxlm-base} \\
    Epoch tối đa & 150 & 20 \\
    Batch size mỗi GPU & 2 & 8 \\
    Gradient accumulation & 8 & 4 \\
    Learning rate & $2\times10^{-5}$ & $2\times10^{-5}$ \\
    LR scheduler & Cosine & Linear \\
    Warm-up ratio & 0,10 & 0,10 \\
    Mixed precision & FP16 & FP16 \\
    Early stopping patience & 15 & 5 \\
    \bottomrule
  \end{tabular}
\end{table}

Các giá trị trong \cref{tab:hyperparams} phản ánh file cấu hình hiện tại. Chúng không chứng minh số epoch thực chạy, checkpoint tốt nhất hay thời gian huấn luyện.

% ============================================================================
\subsection{Phạm vi dữ liệu của từng pipeline}
\label{subsec:training_data_scope}

\begin{table}[htbp]
  \centering
  \caption{Số mẫu có thể sử dụng ở từng giai đoạn.}
  \label{tab:model-data-scope}
  \begin{tabular}{lccc}
    \toprule
    \textbf{Phương pháp} & \textbf{Train} & \textbf{Validation} & \textbf{Test} \\
    \midrule
    Baseline & Không huấn luyện & Không áp dụng & 234 \\
    \model{Donut} & 1.091 & 234 & 234 \\
    \model{LayoutXLM} & 784 & 168 & 234 \\
    \bottomrule
  \end{tabular}
\end{table}

Donut chỉ cần ảnh và nhãn JSON nên sử dụng cả hai nguồn dữ liệu. LayoutXLM cần hộp bao trường để sinh nhãn BIO và do đó chỉ huấn luyện trên các mẫu \code{json\_and\_boxes}. Tất cả phương pháp được chấm trên cùng 234 ID test, nhưng điều kiện huấn luyện khác nhau là một yếu tố nhiễu cần được giữ trong diễn giải.

% ============================================================================
\subsection{Quy trình suy luận và đánh giá}
\label{subsec:evaluation_protocol}

Predictions của ba phương pháp được lưu theo cùng schema, gồm kết quả thô, kết quả chuẩn hóa, trạng thái xử lý và các thành phần latency. Evaluator ghép prediction với ground-truth theo \field{id}. Prediction thiếu hoặc có \code{status=error} được chấm như dự đoán rỗng thay vì bị loại khỏi mẫu số.

Kết quả được báo cáo theo hai chế độ:
\begin{enumerate}
  \item \textbf{All samples:} toàn bộ 234 mẫu; trường hợp cả prediction và ground-truth rỗng được tính đúng theo quy ước metric.
  \item \textbf{Present-only:} chỉ các mẫu có ground-truth khác rỗng của từng trường. Chế độ này được ưu tiên khi phân tích năng lực trích xuất thực sự.
\end{enumerate}

Các tệp JSON metrics, CSV phân tích lỗi, biểu đồ và bảng LaTeX được sinh lại từ predictions hiện có. Bảng trong chương kết quả được nạp từ thư mục \code{report/generated}, tránh sao chép số liệu thủ công.

% ============================================================================
\subsection{Phạm vi phép đo latency}
\label{subsec:latency_protocol}

Artifact prediction chỉ cho phép kết luận về thành phần thời gian đã được ghi:
\begin{itemize}
  \item Baseline: thời gian áp dụng luật và hậu xử lý trên OCR cache.
  \item LayoutXLM: thời gian suy luận mô hình từ OCR cache.
  \item Donut: thời gian suy luận trực tiếp từ ảnh đầu vào (end-to-end từ ảnh sang chuỗi).
\end{itemize}

Thời gian chạy PaddleOCR và VietOCR từ ảnh thô không được ghi đồng nhất trong các prediction hiện có. Vì vậy bảng latency chính chỉ đặt Baseline và LayoutXLM cạnh nhau vì cùng bắt đầu từ OCR cache; báo cáo không cộng một hằng số OCR giả định và không gọi đây là latency end-to-end từ ảnh thô.

% ============================================================================
\subsection{Trạng thái bằng chứng huấn luyện}
\label{subsec:training_evidence}

Các thực nghiệm chạy lại trên Kaggle đã ghi nhận đầy đủ tệp \code{trainer\_state.json} của quá trình huấn luyện:
\begin{itemize}
  \item \textbf{LayoutXLM (ocr\_cache):} Tiến trình dừng tại bước 200 (tương đương 15,4 epochs) với F1 đạt 83,65\% trên tập validation.
  \item \textbf{LayoutXLM (oracle\_ocr):} Chạy đầy đủ 80 bước (20 epochs) trên tập dữ liệu MC-OCR rút gọn có ground-truth OCR (256 mẫu train).
\end{itemize}
Các log huấn luyện chi tiết này hiện được lưu trữ trong thư mục \code{outputs/training\_logs/} làm bằng chứng kiểm chứng và tái lập kết quả.
